\section{Lý do chọn đề tài}

Trong các hệ thống tài chính, ngân hàng và thương mại điện tử hiện đại, mọi thao tác thay đổi dữ liệu --- từ cập nhật trạng thái đơn hàng, điều chỉnh số dư tài khoản, đến sửa thông tin nhân sự --- đều cần được ghi nhận đầy đủ và trung thực. Nhu cầu này xuất phát từ ba áp lực thực tiễn:

\begin{itemize}
    \item \textbf{Tuân thủ pháp lý và kiểm toán:} Các chuẩn mực như \ac{sox}, \ac{pci}, ISO 27001 yêu cầu tổ chức có khả năng trả lời câu hỏi ``ai đã thay đổi gì, vào lúc nào, giá trị cũ và mới là bao nhiêu'' cho bất kỳ bản ghi nào trong vòng nhiều năm. Thiếu \textit{audit trail} đầy đủ dẫn đến không vượt qua kiểm toán và chịu phạt hành chính.
    \item \textbf{Phát hiện gian lận và điều tra sự cố:} Chuỗi thay đổi bất thường --- cùng một tài khoản thực hiện loạt giao dịch nhỏ liên tiếp, hoặc trạng thái đơn hàng bị sửa ngược chiều quy trình --- là dấu hiệu phát hiện sớm gian lận. Không có audit log, điều tra trở nên không thể.
    \item \textbf{Khôi phục và phân tích nguyên nhân gốc rễ:} Khi xảy ra lỗi dữ liệu hoặc bug hệ thống, audit log cho phép \textit{replay} lại chuỗi sự kiện để xác định điểm gãy và phục hồi trạng thái đúng.
\end{itemize}

Tuy nhiên, triển khai hệ thống Audit Log trong môi trường \textbf{write-heavy} (hàng nghìn giao dịch/giây) đặt ra ba thách thức kỹ thuật lớn đồng thời:

\begin{itemize}
    \item \textbf{Thách thức lưu trữ:} Mỗi giao dịch sinh ít nhất một log row với đầy đủ giá trị cũ/mới --- khối lượng tích lũy theo thời gian lên đến hàng triệu dòng, hàng chục GB. Quản lý retention, archiving và indexing trở thành thách thức lớn ở quy mô này.
    \item \textbf{Thách thức hiệu năng:} Mỗi \ac{dml} nghiệp vụ kéo theo ít nhất một thao tác ghi log bổ sung. Nếu overhead này không được kiểm soát, throughput hệ thống suy giảm, ảnh hưởng trực tiếp đến trải nghiệm người dùng và \ac{sla}.
    \item \textbf{Thách thức an toàn:} Bản thân audit log là mục tiêu tấn công --- kẻ gian có thể xóa hoặc sửa log sau khi thực hiện hành vi gian lận để xóa bằng chứng. Hệ thống cần cơ chế \textbf{bất biến} (immutability) và \textbf{chống giả mạo có thể kiểm chứng} (tamper-evident) độc lập với quyền database admin.
\end{itemize}

Đề tài này nghiên cứu và xây dựng giải pháp giải quyết đồng thời ba thách thức trên, thuần túy dựa vào các cơ chế gốc của PostgreSQL 16 (Trigger, JSONB, Declarative Partitioning, SECURITY DEFINER, pgcrypto) mà không phụ thuộc công cụ bên ngoài, phù hợp với hạ tầng cơ sở dữ liệu quan hệ phổ biến hiện nay.

\section{Bài toán nghiên cứu}

\begin{itemize}
    \item \textbf{Input:} Hệ thống PostgreSQL có nhiều bảng nghiệp vụ (đơn hàng, sản phẩm, tài khoản\ldots) đang chịu tải ghi cao; nhiều người dùng với vai trò khác nhau (\texttt{app\_user}, \texttt{db\_admin}) thực hiện các thao tác \ac{dml} (INSERT/UPDATE/DELETE) đồng thời liên tục.
    \item \textbf{Output:} Bảng \texttt{audit\_logs} partitioned JSONB ghi lại đầy đủ mọi thay đổi --- ai thực hiện, thao tác gì, thời điểm nào, giá trị trước và sau --- kèm cơ chế truy vấn nhanh (GIN + partition pruning), bảo vệ bất biến (WORM) chống can thiệp từ mọi cấp quyền, và bằng chứng toán học kiểm chứng tính toàn vẹn (hash chain SHA-256).
\end{itemize}

\section{Mục tiêu nghiên cứu}

\begin{itemize}
    \item \textbf{Mục tiêu tổng quát:} Thiết kế và triển khai kiến trúc Audit Log hiệu năng cao trên PostgreSQL 16, tích hợp trực tiếp vào tầng cơ sở dữ liệu: ghi log tự động không cần sửa mã ứng dụng, lưu trữ linh hoạt đa cấu trúc, truy xuất nhanh theo thời gian và nội dung JSONB, và đảm bảo tính bất biến có thể kiểm chứng độc lập.
    \item \textbf{Mục tiêu cụ thể:}
    \begin{enumerate}
        \item \textbf{Lưu trữ} linh hoạt đa cấu trúc bằng JSONB và truy vấn có cấu trúc.
        \item \textbf{Xử lý/Hiệu năng}: đảm bảo overhead thấp, không làm chậm giao dịch nghiệp vụ.
        \item \textbf{An toàn-bảo mật}: đảm bảo bất biến và chống giả mạo (immutability, tamper-evident).
    \end{enumerate}
\end{itemize}

\section{Câu hỏi nghiên cứu}

\begin{itemize}
    \item \textbf{RQ1:} Kiến trúc Trigger + JSONB + Partitioning có duy trì overhead dưới ngưỡng chấp nhận được ($< 15\%$ TPS) trong điều kiện write-heavy với 50 clients đồng thời không?
    \item \textbf{RQ2:} GIN Index cải thiện hiệu năng truy vấn theo nội dung JSONB đến mức nào, và trong điều kiện nào GIN vượt trội so với Sequential Scan?
    \item \textbf{RQ3:} Cơ chế Append-only/WORM kết hợp hash chain SHA-256 có đủ để ngăn chặn can thiệp log và cung cấp bằng chứng kiểm chứng được không?
\end{itemize}

\section{Đóng góp của đề tài}

\begin{itemize}
    \item (C1) Mô hình audit log schema-less bằng JSONB + GIN.
    \item (C2) Bộ generic trigger/function hỗ trợ audit đa bảng.
    \item (C3) Thiết kế partitioning theo thời gian + retention/archiving.
    \item (C4) Thiết kế bảo mật: security definer, append-only/WORM, tamper-evident.
    \item (C5) Bộ kịch bản benchmark/đánh giá định lượng.
\end{itemize}

\section{Bố cục báo cáo}

Báo cáo được tổ chức thành 6 chương và phần phụ lục:

\begin{itemize}
    \item \textbf{Chương 1 --- Mở đầu:} Trình bày lý do chọn đề tài, bài toán nghiên cứu, mục tiêu, câu hỏi nghiên cứu và đóng góp của đề tài.
    \item \textbf{Chương 2 --- Tổng quan và cơ sở lý thuyết:} Trình bày tổng quan các giải pháp audit log hiện có (WAL-based, Extension, CDC), phân tích lý do lựa chọn hướng tiếp cận Trigger + JSONB + Partitioning trên PostgreSQL.
    \item \textbf{Chương 3 --- Phương pháp đề xuất và thiết kế hệ thống:} Mô tả kiến trúc tổng quan, thiết kế schema bảng \texttt{audit\_logs} (partitioned, JSONB), chiến lược indexing và phân tích đánh đổi kỹ thuật.
    \item \textbf{Chương 4 --- Triển khai cơ chế ghi log, vòng đời và bảo mật:} Chi tiết hàm trigger PL/pgSQL generic, quản lý retention/archiving theo partition, mô hình phân quyền ba lớp, cơ chế Immutability (WORM) + dblink autonomous transaction, tamper-evident hash chain SHA-256.
    \item \textbf{Chương 5 --- Thực nghiệm, kiểm thử và đánh giá:} Kết quả 5 kịch bản benchmark (TPS/latency trên scaling curve 10--80 clients, partitioning/retention, bảo mật, JSONB query performance, DDL audit) với số liệu định lượng và phân tích.
    \item \textbf{Chương 6 --- Kết luận và hướng phát triển:} Trả lời 3 câu hỏi nghiên cứu và đề xuất 4 hướng phát triển tiếp theo.
    \item \textbf{Phụ lục A--D:} DDL đầy đủ, code PL/pgSQL, script benchmark và verify, truy vấn mẫu kiểm toán.
\end{itemize}
